\section{Results}
\label{sec:results}

\subsection{As found: a linearity pass concealing a 424\,\% error}
\label{sec:trap}

The instrument as received declared $\Rcal = \SI{10}{\kilo\ohm}$ in firmware
while carrying the \SI{200}{\ohm} resistor in the \texttt{RCAL1} position. Under
that configuration a cyclic voltammogram of the we-C reference cell was smooth,
symmetric and highly linear, fitting Ohm's law with $R^2 = 0.99999647$. The
recovered resistance was $\SI{55345}{\ohm}$ against a nominal
$\Rnom$---an error of $+424.1\,\%$, equivalently reported currents low by a
factor of $5.22$. We quote the error in recovered resistance throughout, since
that is the quantity compared against the traceable reference; the corresponding
current error is $-80.9\,\%$.

To establish that the declared calibration constant was the cause rather than a
correlate, it was swept deliberately over the board's available values while
every other setting was held fixed, three replicates each
(Table~\ref{tab:rcal}, Fig.~\ref{fig:trap}). The recovered resistance moved from
$\SI{10606}{\ohm}$ to $\SI{55345}{\ohm}$, a factor of $5.22$, tracking the
calibrated transimpedance $\hat{R}_{\mathrm{TIA}}$ exactly: the ratio of
recovered resistances between any two conditions equals the ratio of their
calibrated transimpedance values to within $0.1\,\%$. Restoring the declaration
to \SI{200}{\ohm} returned the instrument to $+0.43\,\%$ reproducibly, including
after passing through the erroneous settings.

Over that sweep the coefficient of determination of the twelve individual fits
lay between $0.99999635$ and $0.99999874$. The absolute error of the condition
means spanned a factor of $982$; $R^2$ varied by $2.4\times10^{-6}$, and not
monotonically. Every condition,
including the $+424\,\%$ one, would pass any linearity-based acceptance
threshold that has appeared in this literature. No feature of the
voltammograms---noise, symmetry, curvature, offset---distinguishes them
(Fig.~\ref{fig:trap}a).

Two details of Table~\ref{tab:rcal} bear on correctability. First, declaring
\SI{1}{\kilo\ohm} and \SI{4.7}{\kilo\ohm} produced nearly the same calibrated
transimpedance ($\SI{2673.9}{\ohm}$ and $\SI{2680.5}{\ohm}$, differing by
$0.25\,\%$) despite a $4.7\times$ difference in the declared value. The gain
error is therefore \emph{not} proportional to the mis-declaration, as
Eq.~(\ref{eq:rtiacal}) alone would predict. This is the saturation anticipated
in Section~\ref{sec:traceability}: the calibration excitation is sized from the
declared value, and once it over-drives the transimpedance output the measured
ratio is set by the clipping level rather than by the resistor. Second, it
follows that the error cannot be removed in post-processing by any factor
derivable from the declared and fitted values. It has to be prevented at the
source.

\begin{table}[!t]
\caption{Effect of the declared calibration-resistor value on recovered
resistance. The fitted resistor is \SI{200}{\ohm} throughout; only the firmware
declaration changes. Three replicate sweeps per condition,
$-200$ to $+\SI{200}{\milli\volt}$, \SI{5}{\milli\volt} steps,
\SI{200}{\milli\volt\per\second}.}
\label{tab:rcal}
\centering
\footnotesize
\begin{tabular}{r r r r r}
\toprule
\textbf{Declared} & \textbf{Calibrated} & \textbf{Recovered} &
\textbf{Error vs.} & \textbf{Mean} \\
$\Rcal$ ($\Omega$) & $\hat{R}_{\mathrm{TIA}}$ ($\Omega$) & $R$ ($\Omega$) &
nominal (\%) & $R^{2}$ \\
\midrule
\textbf{200}$^{*}$ & 1093.0 & $10\,605.6 \pm 0.4$  & $+0.43$   & 0.99999861 \\
1\,000                & 2673.9 & $25\,944.4 \pm 4.8$  & $+145.7$  & 0.99999831 \\
4\,700                & 2680.5 & $26\,010.3 \pm 4.6$  & $+146.3$  & 0.99999831 \\
10\,000               & 5703.1 & $55\,345.5 \pm 11.9$ & $+424.1$  & 0.99999647 \\
\midrule
\multicolumn{3}{l}{\emph{range across the series}} & $\times 982$ & $2.1\times10^{-6}$ \\
\bottomrule
\end{tabular}
\\[2pt]
\parbox{\linewidth}{\footnotesize $^{*}$ The correct declaration, matching the
fitted resistor. Uncertainties are standard deviations over three replicates.}
\end{table}

\begin{figure*}[!t]
\centering
\includegraphics[width=0.92\textwidth]{fig2_linearity_trap}
\caption{Linearity is blind to gain error. (a) Voltammograms of the same
reference cell under four declared calibration-resistor values; the dashed line
is the response computed from component nominals. Only the \SI{200}{\ohm}
declaration overlies it. The \SI{1}{\kilo\ohm} trace falls on the
\SI{4.7}{\kilo\ohm} trace because the calibration excitation has saturated.
(b) Absolute error and unexplained variance $100(1-R^{2})$ on the same
logarithmic axis. The error spans three decades while $R^{2}$ is flat to
$\sim\!10^{-6}$.}
\label{fig:trap}
\end{figure*}

\subsection{Firmware defects identified}
\label{sec:defects}

Diagnosis of the scale error, and a systematic comparison of the failing method
classes against the working sequencer-driven voltammetry engine and against the
vendor reference examples, produced the eight defects in
Table~\ref{tab:defects}. All were corrected; the corrections are the firmware
state used for every subsequent measurement reported here.

Three points are worth drawing out. The mis-declared constant (D1) was, on its
own, correctable at run time---except that the user-facing command intended to
set it wrote a parameter store that no measurement code reads, while the
voltammetry and impedance paths read a separate global written only by an
unreachable legacy handler (D2). The instrument therefore reported one
calibration value through its parameter listing and used another, and the
mismatch could not be fixed without recompiling. Even after D2 was repaired, the
transimpedance calibration ran only once per power-up (D3), so a corrected
declaration would not have taken effect until the next reset. Three separately
minor faults thus combined into an error that was neither visible nor
correctable from the host.

\begin{table*}[!t]
\caption{Firmware defects identified and corrected in this work.}
\label{tab:defects}
\centering
\footnotesize
\setlength{\tabcolsep}{4pt}
\begin{tabular}{@{}l >{\raggedright\arraybackslash}p{3.1cm} >{\raggedright\arraybackslash}p{4.2cm} >{\raggedright\arraybackslash}p{4.0cm}@{}}
\toprule
\textbf{ID} & \textbf{Site} & \textbf{Defect} & \textbf{Observed effect} \\
\midrule
D1 & \texttt{AD5941\_25.ino} (global \texttt{fRcal}) & Declared $\Rcal$ = \SI{10}{\kilo\ohm}; fitted resistor is \SI{200}{\ohm} & Currents low by $5.22\times$, i.e.\ $+424\,\%$ in recovered resistance; invisible to $R^{2}$ \\
D2 & \texttt{communication.cpp} (\texttt{c} command, parameter listing) & Command writes a parameter store no measurement path reads; voltammetry and impedance read a separate global written only by unreachable legacy code & $\Rcal$ not settable at run time; reported value $\neq$ value in use \\
D3 & \texttt{rampTest.cpp} (\texttt{AppRAMPInit}) & Transimpedance calibration gated on a flag that latches after the first sweep & Calibration frozen at power-up; later corrections ignored \\
D4 & \texttt{utilities.cpp} (predicate \texttt{Delay}) & Unbounded busy-wait on a hardware flag, polled at \SI{100}{\milli\second} & Impedance sweep hangs indefinitely; instrument unresponsive while USB still enumerates \\
D5 & \texttt{c\_ca/c\_swv/c\_dpv.cpp} & \texttt{ADCPga = 1} selects \texttt{ADCPGA\_1P5} (gain 1.5) while the conversion divides by 1.0 & Silent $1.5\times$ current error in all three DC methods \\
D6 & \texttt{c\_ca/c\_swv/c\_dpv.cpp} & \texttt{LPDACSW\_VZERO2PIN} closed, tying the $V_{\mathrm{zero}}$ buffer to the sense pin & Cell current sunk by the buffer instead of the transimpedance stage \\
D7 & \texttt{c\_ca/c\_swv/c\_dpv.cpp} & Bias code added to rather than subtracted from the $V_{\mathrm{zero}}$ offset & Applied-potential polarity inverted \\
D8 & \texttt{c\_ca/c\_swv/c\_dpv.cpp} & AFE never woken; the voltammetry path leaves it hibernating with the low-power loop off & Register writes succeed while the analog front end stays powered down \\
\midrule
D9 & \texttt{cv.cpp} (\texttt{\_AD5940\_Main}) & Unbounded wait on the sequencer completion flag & Instrument wedged partway through an extended sweep campaign and required a power cycle \\
\bottomrule
\end{tabular}
\\[2pt]
\parbox{\textwidth}{\footnotesize D5--D8 existed as three verbatim copies each,
one per method class, and a fourth unreachable copy in a legacy free-function
module. D9 was encountered during the measurement campaign reported here; after
bounding the wait, the same condition reports \texttt{ERR,CV\_TIMEOUT} and the
instrument remains responsive, though a power cycle is still required to restore
full performance (Section~\ref{sec:invariance}).}
\end{table*}

\subsection{After correction: accuracy and repeatability}
\label{sec:corrected}

With the declaration corrected, ten replicate sweeps under identical settings
recovered
$R = (10\,605.2 \pm 1.6)~\Omega$
against the nominal $\Rnom$, an error of $+0.428\,\%$ with individual replicates
spanning $+0.403$ to $+0.452\,\%$ (Fig.~\ref{fig:repeat}). The relative standard
deviation was $0.0152\,\%$, some thirty times inside the $0.5\,\%$ acceptance
window of check V4, and the accuracy is well inside the $2\,\%$ window of V2.

The calibrated transimpedance was $(1093.02 \pm 0.03)~\Omega$ across the same
ten sweeps, a relative standard deviation of $0.0023\,\%$. The front end
therefore measures its own gain to a few parts in $10^{5}$; the as-found defect
was not a noisy calibration but a precise measurement against the wrong
reference.

Check V3 is informative about what limits the fit. The median root-mean-square
residual was \SI{12.6}{\nano\ampere}, against a current quantisation step of
$\mathrm{LSB} = \SI{50.8}{\nano\ampere}$ at the calibrated gain---that is,
$0.25$~LSB. The residuals show no structure against potential
(Fig.~\ref{fig:corrected}b). The voltammetry fit is thus limited by the
converter, not by the analog chain, and $R^2$ has no further headroom to report:
this is the regime in which a linearity statistic is least informative and most
reassuring.

\begin{figure}[!t]
\centering
\includegraphics[width=0.92\linewidth]{fig3_corrected_fit}
\caption{Corrected instrument against the reference cell. (a) A single sweep
with the response computed from component nominals overlaid. (b) Fit residuals,
which are unstructured and of order a quarter of the current quantisation
step.}
\label{fig:corrected}
\end{figure}

\begin{figure}[!t]
\centering
\includegraphics[width=0.92\linewidth]{fig4_repeatability}
\caption{Ten replicate sweeps under identical settings. The shaded band is
$\pm1$ standard deviation; the dashed line is the value computed from component
nominals. The offset from nominal is systematic, not scatter.}
\label{fig:repeat}
\end{figure}

\subsection{Invariance across acquisition settings, and an operating limit}
\label{sec:invariance}

Check V5 requires the recovered resistance to be a property of the cell rather
than of the settings. Across scan rates from \SIrange{25}{800}{\milli\volt\per\second},
potential limits from $\pm100$ to $\pm\SI{600}{\milli\volt}$, and potential
steps from \SIrange{5}{20}{\milli\volt}, the recovered resistance spanned
$+0.411\,\%$ to $+0.485\,\%$ error---a total spread of $0.074$ percentage
points, well inside the acceptance window, with $R^2 \ge 0.999994$ throughout
(Fig.~\ref{fig:robust}a--c). The residual offset is common to every
configuration, which is the signature of a scale term rather than a
settings-dependent artefact. One further step size, \SI{2}{\milli\volt}, was
attempted and failed; it falls below the timing limit established next rather
than indicating a resolution boundary.

One boundary emerged. Sweeps failed outright---returning no data, and in the
more extreme cases the bounded-wait timeout---whenever the potential step period
$E_{\mathrm{step}}/\nu$ fell below about \SI{19}{\milli\second}
(Fig.~\ref{fig:robust}d). The threshold lies between the \SI{20}{\milli\second}
condition, which succeeded, and the \SI{17.5}{\milli\second} condition, which did
not. This is roughly $2.7\times$ the \SI{7}{\milli\second} settling delay the
firmware configures, so the achievable sweep speed cannot be read off that
parameter. The limit is on the step \emph{period}, not on scan rate as such:
\SI{800}{\milli\volt\per\second} is accurate to $+0.435\,\%$ when taken in
\SI{20}{\milli\volt} steps, while \SI{400}{\milli\volt\per\second} fails in
\SI{5}{\milli\volt} steps. At \SI{5}{\milli\volt} resolution the usable maximum
is \SI{250}{\milli\volt\per\second}.

Bounding the acquisition wait (D9) prevents the instrument from wedging, but it
does not fully restore the front end. Sweeps taken immediately after a timeout
were reproducibly degraded---$R^{2}$ fell from $0.999999$ to about $0.9999$ at
unchanged accuracy---and recovered only after a power cycle, which we confirmed
twice. Reported measurements were therefore all acquired on a freshly reset
instrument, and we recommend that a timeout be treated as requiring a reset
rather than as a recoverable condition.

\begin{figure*}[!t]
\centering
\includegraphics[width=0.95\textwidth]{fig5_robustness}
\caption{Parameter invariance and the operating limit. (a)--(c) Recovered
resistance error against scan rate, potential limit and potential step; the
dashed line is the ten-replicate mean of Fig.~\ref{fig:repeat}. The total spread
is $0.074$ percentage points. (d) Sweep outcome against potential step period:
acquisition fails below about \SI{19}{\milli\second}, roughly $2.7\times$ the
settling delay configured in firmware.}
\label{fig:robust}
\end{figure*}

\subsection{The capacitive term is below the resolution floor}
\label{sec:cap}

Equation~(\ref{eq:ramp}) predicts that the forward and reverse branches separate
by $\Delta I = 2\nu C_{\mathrm{eff}}$, which for
$C_{\mathrm{eff}} = \SI{29.6}{\nano\farad}$ gives \SI{1.5}{\nano\ampere} at
\SI{25}{\milli\volt\per\second} rising to \SI{47.4}{\nano\ampere} at
\SI{800}{\milli\volt\per\second}. The current quantisation step is
\SI{50.8}{\nano\ampere}. The predicted separation is therefore below one LSB at
every scan rate the instrument supports, reaching at most $0.93$~LSB.

Measured branch separations were consistent with this: across the scan-rate
series they scattered between $-2$ and $+\SI{7}{\nano\ampere}$ with no
dependence on scan rate, where a resolvable capacitance would have produced a
proportional trend spanning \SIrange{1.5}{47}{\nano\ampere}. We therefore report
that the voltammetry path constrains only the resistive part of the reference
cell, and that recovering $C_{dl}$ from it is outside the instrument's present
resolution rather than merely unattempted.

\subsection{Method coverage}
\label{sec:coverage}

The same reference cell, board and session were used to exercise the remaining
implemented methods (Table~\ref{tab:coverage}). Cyclic voltammetry is the only
method for which we can report quantitative agreement.

Chronoamperometry, square-wave and differential pulse voltammetry share a
direct-register implementation of the low-power loop. After correcting D5--D8,
the analog-to-digital converter demonstrably converts---every read completed
without timeout, and a genuine transient follows each potential step---but the
steady-state current does not depend on the applied potential. A step of
$+\SI{200}{\milli\volt}$ across the \SI{10.56}{\kilo\ohm} reference should draw
\SI{18.9}{\micro\ampere}; the measured steady state was pinned near the
quantisation floor and, after the corrections, at a fixed offset, with the
difference between $\pm\SI{200}{\milli\volt}$ steps below \SI{3}{\nano\ampere}
against an expected \SI{37.9}{\micro\ampere}. Sweeping the transimpedance output
topology over four representative filter-resistor settings and both output
switch configurations---eight combinations, spanning a bypassed filter to the
\SI{1}{\mega\ohm} setting---changed nothing, which places the fault upstream of
the sensing stage, in the
establishment of the cell drive, rather than in the conversion or scaling path
that the voltammetry engine shares.

Impedance spectroscopy did not complete a single frequency point. In a sweep
configured for fifteen logarithmically spaced frequencies between
\SI{100}{\hertz} and \SI{100}{\kilo\hertz}, the discrete Fourier transform ready
flag failed to assert at every point the acquisition reached, and the sweep
yielded no valid impedance value. Before correction
of D4 this manifested as an indefinite hang; afterwards the sweep reports
timeouts and returns control in \SI{50}{\second}. Open-circuit potential
returned $-\SI{1820}{\milli\volt}$, exactly the negative of the converter
reference, i.e.\ a saturated reading.

\begin{table}[!t]
\caption{Method coverage on the we-C reference cell. Only cyclic voltammetry is
claimed as validated.}
\label{tab:coverage}
\centering
\footnotesize
\begin{tabular}{l l >{\raggedright\arraybackslash}p{3.6cm}}
\toprule
\textbf{Method} & \textbf{Status} & \textbf{Evidence} \\
\midrule
CV  & \textbf{Validated} & $+0.428\,\%$, \%RSD $0.0152$, invariant across settings \\
CA  & Not functional & Steady-state current independent of applied potential \\
SWV & Not functional & Shares the defective drive path with CA \\
DPV & Not functional & Shares the defective drive path with CA \\
EIS & Not functional & DFT ready flag never asserts; now fails in \SI{50}{\second} rather than hanging \\
OCP & Not functional & Returns $-V_{\mathrm{ref}}$ (saturated) \\
\bottomrule
\end{tabular}
\end{table}
